\chapter{\texorpdfstring{$F$}{F}分布}
    \section{定義}
        \begin{shadebox}
            関数$f(f)$を次のように定義する。
            \[f(f) = \frac{\Gamma{\left(\frac{\nu_1+\nu_2}{2}\right)}}{\Gamma{\left(\frac{\nu_1}{2}\right)}\Gamma{\left(\frac{\nu_2}{2}\right)}}\left(\frac{\nu_1}{\nu_2}\right)^{\frac{\nu_1}{2}}f^{\frac{\nu_1-2}{2}}\left(1+\frac{\nu_1}{\nu_2}f\right)^{-\frac{\nu_1+\nu_2}{2}}\]
            \[\text{where} \hspace{10pt} \nu_1,\nu_2 \in \mathbb{N},\;\;\; f > 0\]
            確率密度関数が上式で与えられるような分布を「自由度$(\nu_1,\nu_2)$の$F$分布」と呼び、$F_{\nu_2}^{\nu_1}$と書くことにする。
        \end{shadebox}
    \section{\texorpdfstring{$X_1,X_2 : \text{indep},\;\;\;X_1 \sim {\chi^2}_{\nu_1},\;\;\;X_2 \sim {\chi^2}_{\nu_2} \hspace{10pt} \Rightarrow \hspace{10pt} F\coloneqq \frac{\frac{X_1}{\nu_1}}{\frac{X_2}{\nu_2}} \sim F_{\nu_2}^{\nu_1}$}{X1,X2: 独立, X1～χ\string^2\_ν1, X2～χ\string^2\_ν2 ⇒ F := (X1/ν1)/(X2/ν2) ～ F\_ν2\string^ν1}}
        \begin{shadebox}
            \[X_1,X_2 : \text{indep},\;\;\;X_1 \sim {\chi^2}_{\nu_1},\;\;\;X_2 \sim {\chi^2}_{\nu_2} \hspace{10pt} \Rightarrow \hspace{10pt} F\coloneqq \frac{\frac{X_1}{\nu_1}}{\frac{X_2}{\nu_2}} \sim F_{\nu_2}^{\nu_1}\]
        \end{shadebox}
        \begin{proof}
            \quad\par
            $X_1,X_2,F$の確率密度関数を$f^{X_1}(x_1),f^{X_2}(x_2),f(f)$と書くことにする。
            \begin{align*}
                P(0 \leq F \leq f) &= P\left(0 \leq X_1 \leq \frac{\nu_1}{\nu_2}x_2f\right) = \int_{x_2=0}^{\infty}{\int_{x_1=0}^{\frac{\nu_1}{\nu_2}x_2f}{f^{X_1}(x_1)f^{X_2}(x_2) dx_1} dx_2}\\
                &= \int_{x_2=0}^{\infty}{f^{X_2}(x_2)\left(\int_{x_1=0}^{\frac{\nu_1}{\nu_2}x_2f}{f^{X_1}(x_1) dx_2}\right) dx_1}
            \end{align*}
            \begin{align}
                \therefore \;\;\; f(f) &= \frac{\partial}{\partial f}P(0 \leq F \leq f) = \int_{x_2=0}^{\infty}{f^{X_2}(x_2)f^{X_1}\left(\frac{\nu_1}{\nu_2}x_2f\right)\frac{\nu_1}{\nu_2}x_2 dx_2} \quad \text{(微分と積分の順序を交換した)} \nonumber \\
                &= \frac{1}{2^{\frac{\nu_2}{2}}\Gamma{\left(\frac{\nu_2}{2}\right)}2^{\frac{\nu_1}{2}}\Gamma{\left(\frac{\nu_1}{2}\right)}}\frac{\nu_1}{\nu_2}\int_{0}^{\infty}{x^{\frac{\nu_2}{2}-1}e^{-\frac{x}{2}}\left(\frac{\nu_1}{\nu_2}f\right)^{\frac{\nu_1}{2}-1}x^{\frac{\nu_1}{2}-1}\exp\left[-\frac{x}{2}\left(1+\frac{\nu_1}{\nu_2}f\right)\right] dx} \nonumber \\
                &=\frac{1}{2^{\frac{\nu_1+\nu_2}{2}}\Gamma{\left(\frac{\nu_2}{2}\right)}\Gamma{\left(\frac{\nu_1}{2}\right)}}\left(\frac{\nu_1}{\nu_2}\right)^\frac{\nu_1}{2}f^{\frac{\nu_1-2}{2}}\int_{0}^{\infty}{x^{\frac{\nu_1+\nu_2}{2}-1}\exp\left[-\frac{x}{2}\left(1+\frac{\nu_1}{\nu_2}f\right)\right] dx}
            \end{align}
            先に積分の値を計算する。$\displaystyle \frac{x}{2}\left(1+\frac{\nu_1}{\nu_2}f\right) = u$なる変数変換を行って
            \begin{align*}
                & \int_{0}^{\infty}{x^{\frac{\nu_1+\nu_2}{2}-1}\exp\left[-\frac{x}{2}\left(1+\frac{\nu_1}{\nu_2}f\right)\right] dx} = \int_{0}^{\infty}{\left(\frac{2}{1+\frac{\nu_1}{\nu_2}f}u\right)^{\frac{\nu_1+\nu_2}{2}-1}e^{-u}\frac{2}{1+\frac{\nu_1}{\nu_2}f} du}\\
                &= \left(1+\frac{\nu_1}{\nu_2}f\right)^{-\frac{\nu_1+\nu_2}{2}}2^{\frac{\nu_1+\nu_2}{2}}\int_{0}^{\infty}{u^{\frac{\nu_1+\nu_2}{2}-1}e^{-u} du} = \left(1+\frac{\nu_1}{\nu_2}f\right)^{-\frac{\nu_1+\nu_2}{2}}2^{\frac{\nu_1+\nu_2}{2}}\Gamma{\left(\frac{\nu_1+\nu_2}{2}\right)}
            \end{align*}
            これを(2)に適用して
            \[f(f) = \frac{\Gamma{\left(\frac{\nu_1+\nu_2}{2}\right)}}{\Gamma{\left(\frac{\nu_1}{2}\right)}\Gamma{\left(\frac{\nu_2}{2}\right)}}\left(\frac{\nu_1}{\nu_2}\right)^{\frac{\nu_1}{2}}f^{\frac{\nu_1-2}{2}}\left(1+\frac{\nu_1}{\nu_2}f\right)^{-\frac{\nu_1+\nu_2}{2}} = \text{pdf of} \quad F_{\nu_2}^{\nu_1}\]
        \end{proof}
    \section{自由度\texorpdfstring{$(\nu_1,\nu_2)$}{(ν1,ν2)}の\texorpdfstring{$F$}{F}分布\texorpdfstring{$F_{\nu_2}^{\nu_1}$}{F\_ν2\string^ν1}は\texorpdfstring{$\nu_2 \geq 3$}{ν2≧3}の時に限り期待値が定義できてその値は\texorpdfstring{$\frac{\nu_2}{\nu_2-2}$}{ν2/(ν2-2)}である。}
        \begin{shadebox}
            自由度$(\nu_1,\nu_2)$の$F$分布$F_{\nu_2}^{\nu_1}$は$\nu_2 \geq 3$の時に限り期待値が定義できてその値は$\frac{\nu_2}{\nu_2-2}$である。
        \end{shadebox}
        \begin{proof}
            \setcounter{equation}{0}
            \quad\par
            $\nu_2 \leq 2$の場合、$x\to \infty$で$xf(x)$が$\frac{1}{x}$オーダーを下回らないので$\int_{0}^{\infty}{xf(x) dx}$は発散し期待値は定義できない。
            $\nu_2 \geq 3$の場合、$x\to \infty$で$xf(x)$が$\frac{1}{x}$オーダーを下回るので積分は収束して期待値が定義できる。この値が$\frac{\nu_2}{\nu_2-2}$であることを示す。\\
            \begin{equation}
                E[F] = \int_{0}^{\infty}{xf(x) dx} = \frac{\Gamma{\left(\frac{\nu_1+\nu_2}{2}\right)}}{\Gamma{\left(\frac{\nu_1}{2}\right)}\Gamma{\left(\frac{\nu_2}{2}\right)}}\left(\frac{\nu_1}{\nu_2}\right)^{\frac{\nu_1}{2}}\int_{0}^{\infty}{x^{\frac{\nu_2}{2}}\left(1+\frac{\nu_1}{\nu_2}x\right)^{-\frac{\nu_1+\nu_2}{2}} dx}
            \end{equation}
            積分の値を先に計算する。$\displaystyle u = \frac{\frac{\nu_1}{\nu_2}x}{1+\frac{\nu_1}{\nu_2}x} = \frac{\nu_1x}{\nu_2+\nu_1x}$なる変数変換を行うと$x = \frac{\nu_2}{\nu_1}\frac{u}{1-u}, \quad dx = \frac{\nu_2}{\nu_1}\frac{1}{(1-u)^2}du$となり積分範囲は$u=0\to \textcolor{red}{1}$になる。\\
            (※(1)で積分の左側に括りだした定数群にガンマ関数が居る。最終的にこれらが消えてほしいわけだから積分結果がガンマ関数あるいはベータ関数にならなければならない。そこでベータ関数の出現を期待して積分範囲を$0\to\infty$から$0\to 1$に変更するような変数変換を行ったのである。)\\
            (1)の積分は
            \begin{align*}
                & \int_{0}^{1}{\left(\frac{\nu_2}{\nu_1}\right)^{\frac{\nu_1}{2}}\left(\frac{u}{1-u}\right)^{\frac{\nu_1}{2}}\left(1+\frac{u}{1-u}\right)^{-\frac{\nu_1+\nu_2}{2}}\frac{\nu_2}{\nu_1}\frac{1}{(1-u)^2} du}\\
                &= \left(\frac{\nu_1}{\nu_2}\right)^{-\frac{\nu_1}{2}-1}\int_{0}^{1}{u^{\frac{\nu_1}{2}}(1-u)^{\frac{\nu_2}{2}-2} du} = \left(\frac{\nu_1}{\nu_2}\right)^{-\frac{\nu_1}{2}-1}\int_{0}^{1}{u^{\frac{\nu_1}{2}+1-1}(1-u)^{\frac{\nu_2}{2}-1-1} du}\\
                &= \left(\frac{\nu_1}{\nu_2}\right)^{-\frac{\nu_1}{2}-1}B\left(\frac{\nu_1}{2}+1,\frac{\nu_2}{2}-1\right) = \left(\frac{\nu_1}{\nu_2}\right)^{-\frac{\nu_1}{2}-1}\frac{\Gamma{\left(\frac{\nu_1}{2}+1\right)}\Gamma{\left(\frac{\nu_2}{2}-1\right)}}{\Gamma{\left(\frac{\nu_1+\nu_2}{2}\right)}}\\
                &= \left(\frac{\nu_1}{\nu_2}\right)^{-\frac{\nu_1}{2}-1}\frac{\frac{\nu_1}{2}\Gamma{\left(\frac{\nu_1}{2}\right)}\Gamma{\left(\frac{\nu_2}{2}\right)}}{\left(\frac{\nu_2}{2}-1\right)\Gamma{\left(\frac{\nu_1+\nu_2}{2}\right)}}
            \end{align*}
            これを(1)に適用して
            \[E[F] = \frac{\nu_2}{\nu_2-2}\]
        \end{proof}
    \section{自由度\texorpdfstring{$(\nu_1,\nu_2)$}{(ν1,ν2)}の\texorpdfstring{$F$}{F}分布\texorpdfstring{$F_{\nu_2}^{\nu_1}$}{F\_ν2\string^ν1}の分散は\texorpdfstring{$\nu_2 \leq 4$}{ν2<4}の場合\texorpdfstring{$\infty$}{∞}、\texorpdfstring{$\nu_2 \geq 5$}{ν2≧5}の場合\texorpdfstring{$\displaystyle \frac{2{\nu_2}^2(\nu_1+\nu_2-2)}{\nu_1(\nu_2-2)^2(\nu_2-4)}$}{(2(ν2)\string^2 (ν1+ν2-2))/(ν1(ν2-2)\string^2 (ν2-4))}である。}
        \begin{shadebox}
            自由度$(\nu_1,\nu_2)$の$F$分布$F_{\nu_2}^{\nu_1}$の分散は$\nu_2 \leq 4$の場合$\infty$、$\nu_2 \geq 5$の場合$\displaystyle \frac{2{\nu_2}^2(\nu_1+\nu_2-2)}{\nu_1(\nu_2-2)^2(\nu_2-4)}$である。
        \end{shadebox}
        \begin{proof}
            \setcounter{equation}{0}
            \quad\par
            $\nu_2 \leq 4$の場合、$x\to \infty$で$x^2f(x)$が$\frac{1}{x}$オーダーを下回らないので$\int_{0}^{\infty}{x^2f(x) dx} = \infty$\\
            $\nu_2 \geq 5$の場合、$x\to \infty$で$x^2f(x)$が$\frac{1}{x}$オーダーを下回るので積分は収束する。このとき分散が$\displaystyle \frac{2{\nu_2}^2(\nu_1+\nu_2-2)}{\nu_1(\nu_2-2)^2(\nu_2-4)}$であることを示す。\\
            $V[F] = E[F^2] - {E[F]}^2$であり第2項は定理2より既知である。第1項を計算する。
            \begin{equation}
                E[F^2] = \int_{0}^{\infty}{x^2f(x) dx} = \frac{\Gamma{\left(\frac{\nu_1+\nu_2}{2}\right)}}{\Gamma{\left(\frac{\nu_1}{2}\right)}\Gamma{\left(\frac{\nu_2}{2}\right)}}\left(\frac{\nu_1}{\nu_2}\right)^{\frac{\nu_1}{2}}\int_{0}^{\infty}{x^{\frac{\nu_2}{2}+1}\left(1+\frac{\nu_1}{\nu_2}x\right)^{-\frac{\nu_1+\nu_2}{2}} dx}
            \end{equation}
            積分の値を先に計算する。$\displaystyle u = \frac{\frac{\nu_1}{\nu_2}x}{1+\frac{\nu_1}{\nu_2}x} = \frac{\nu_1x}{\nu_2+\nu_1x}$なる変数変換を行うと$x = \frac{\nu_2}{\nu_1}\frac{u}{1-u}, \quad dx = \frac{\nu_2}{\nu_1}\frac{1}{(1-u)^2}du$となり積分範囲は$u=0\to \textcolor{red}{1}$になる。\\
            (1)の積分は
            \begin{align*}
                & \int_{0}^{1}{\left(\frac{\nu_2}{\nu_1}\right)^{\frac{\nu_1}{2}+1}\left(\frac{u}{1-u}\right)^{\frac{\nu_1}{2}+1}\left(1+\frac{u}{1-u}\right)^{-\frac{\nu_1+\nu_2}{2}}\frac{\nu_2}{\nu_1}\frac{1}{(1-u)^2} du}\\
                &= \left(\frac{\nu_1}{\nu_2}\right)^{-\frac{\nu_1}{2}-2}\int_{0}^{1}{u^{\frac{\nu_1}{2}+1}(1-u)^{\frac{\nu_2}{2}-3} du} = \left(\frac{\nu_1}{\nu_2}\right)^{-\frac{\nu_1}{2}-2}B{\left(\frac{\nu_1}{2}+2,\frac{\nu_2}{2}-2\right)}\\
                &= \left(\frac{\nu_1}{\nu_2}\right)^{-\frac{\nu_1}{2}-2}\frac{\Gamma{\left(\frac{\nu_1}{2}+2\right)}\Gamma{\left(\frac{\nu_2}{2}-2\right)}}{\Gamma{\left(\frac{\nu_1+\nu_2}{2}\right)}}\\
                &= \left(\frac{\nu_1}{\nu_2}\right)^{-\frac{\nu_1}{2}-2}\frac{\left(\frac{\nu_1}{2}+1\right)\frac{\nu_1}{2}\Gamma{\left(\frac{\nu_1}{2}\right)}}{\Gamma{\left(\frac{\nu_1+\nu_2}{2}\right)}}\frac{4\Gamma{\left(\frac{\nu_2}{2}\right)}}{(\nu_2-4)(\nu_2-2)}
            \end{align*}
            これを(1)に適用して
            \[E[F^2] = \frac{{\nu_2}^2(\nu_1+2)}{\nu_1(\nu_2-4)(\nu_2-2)}\]
            定理2で示された$E[F]$を用いて
            \[V[F] = E[F^2]-{E[F]}^2 = \frac{{\nu_2}^2(\nu_1+2)}{\nu_1(\nu_2-4)(\nu_2-2)} - \frac{{\nu_2}^2}{(\nu_2-2)^2} = \frac{2{\nu_2}^2(\nu_1+\nu_2-2)}{\nu_1(\nu_2-2)^2(\nu_2-4)}\]
        \end{proof}